\scanpage{85}
\[
\tag{113}
S_A\Pi M_G\overset{\mathrm{def}}=\Pi M^*_{G,G'},
\]
\margin{ici $A$ désigne l'ens. de toutes les arêtes de $G$.}
\unsure{$\ast$} c'est donc une extension de $\Pi M_{G,G'}$ \margin{$=$\deleted{$(\Pi(M_G,\unsure{\Gamma_G}))$}} par $\mathbf Z^{A(G)}$ \deleted{de}, où $A(G)$ est l'ens. des arêtes de $G$.

Bien entendu, le groupe $\Gamma=\operatorname{Aut}G$ opère non seulement sur $M_G$ et sur $\Pi M_G$, mais aussi sur $S_A\Pi M_G$ \deleted{$=\Gamma$} (via son action sur $M^*_{G,G'}$), de façon à commuter au morphisme canonique
\[
\tag{114}
\diagram{\begin{tikzcd}[ampersand replacement=\&]
S_A\Pi M_G \arrow[r] \& (\Pi M_G\deleted{,\Gamma_G})
\end{tikzcd}}
\]
ce qui permet, pour $\forall$ sous-groupe \deleted{$\Gamma'$}
\[
\Gamma'\subset\Gamma_G
\]
de définir une extension
\margin{expliciter un peu\ldots}
\[
\tag{115}
\diagram{\begin{tikzcd}[ampersand replacement=\&, column sep=large,row sep=small]
(S_A\Pi M_G,\Gamma') \arrow[r] \& (\Pi M_G,\deleted{\unsure{\Gamma_G}}\inserted{\Gamma'})\\
\text{ou } S_A\Pi(M_G,\Gamma') \arrow[r,dash] \& \Pi(M_G,\Gamma')
\end{tikzcd}}
\]
\deleted{Soit $C$ un ensemble d}
dont le « noyau » est le système local $\mathbf Z^{A(G)}$ sur $(\Pi M_G,\Gamma')$, tordu grâce à l'action de $\Gamma'$ sur $A(G)$. Si $C\subset A(G)$ est stable par l'action de $\Gamma'$, alors on a un hom
\[
\tag{116}
\diagram{\begin{tikzcd}[ampersand replacement=\&]
\mathbf Z^{A(G)} \arrow[r] \& \mathbf Z^C
\end{tikzcd}}
\]
\margin{84}
